\documentclass[journal]{IEEEtran}
\usepackage[utf8]{inputenc}
\usepackage[T1]{fontenc}
\usepackage{amsmath,booktabs,url}
\title{AirGuard India: Station-Calibrated Day-Ahead Air Quality Forecasting with Uncertainty-Aware Vulnerability Advisories}
\author{Nishanth R\thanks{Nishanth R is with the Department of AI and Data Science,
Rajalakshmi Engineering College, Chennai, Tamil Nadu, India.
Email: nishanth.r.2024.aids@rajalakshmi.edu.in; backup: nishanth.artwork@gmail.com.}
\thanks{Development manuscript. Primary CAMS results and model selection are pending;
the final evaluation interval remains sealed.}}
\begin{document}
\maketitle
\begin{abstract}
This development manuscript documents an India-specific station-calibration
protocol locked before model fitting; the manuscript itself is not a new
preregistration. The intended comparison is raw operational CAMS versus direct learning
and residual correction, followed by uncertainty calibration and exposure
advisory evaluation. Hourly quality rules, chronological and spatial validation,
and a sealed final test constrain the analysis. The genuine historical CAMS issue-time archive has passed operational GRIB validation; development acquisition is incomplete. A separate preliminary measured-data
study is reported as infrastructure evidence, without a CAMS improvement claim.
Real prospective forecasts are frozen, but matched measured outcomes remain
pending. [RESULT PENDING CAMS ARCHIVE EXPERIMENT]
\end{abstract}
\begin{IEEEkeywords}CAMS, air quality, India, calibration, prospective verification\end{IEEEkeywords}
\section{Problem and Related Work}
The research question is whether a correction learned at Indian stations
improves next-day CPCB-style AQI while retaining useful uncertainty and warning
behaviour. Atmospheric forecasts and measured concentrations have different
provenance. Existing Indian forecasting systems establish that this is an
evaluation problem, not an unoccupied application domain \cite{jena2021delhi}.
Station MOS, CAMS statistical post-processing, conformal intervals and risk
calibration are established methods. The targeted review records overlapping
European CAMS work and distinguishes Indian reanalysis mapping from genuine
issue-time prediction. No world's-first claim is made. The contribution under
test is the audited Indian operational procedure and its empirical warning
trade-offs. Full-text novelty review remains incomplete.
\section{Study Status}
Two streams are kept separate: frozen current forecasts constitute the
prospective study; measured historical targets support non-CAMS preliminary
research. Neither substitutes for the blocked primary archive comparison.








% BEGIN MOS RELATED WORK
Station-specific CAMS post-processing is established. Riccio and Chianese
study statistical correction with spatial verification \cite{riccio2024mos}.
Petetin et al. explicitly distinguish improvements in continuous forecast
errors from categorical ozone-warning skill \cite{petetin2022mos}.
Operational regional CAMS production also includes MOS \cite{colette2025cams}.
The proposed contribution is an audited Indian evaluation joining issue-time
availability, censored-Severe diagnostics and prospective station verification;
neither a new MOS invention nor a world's-first application is claimed.
% END MOS RELATED WORK
% BEGIN ARCHITECTURE
\begin{figure}[h]\centering
\fbox{\begin{minipage}{.91\columnwidth}\centering\small
\textbf{Measured stations}\ $\longrightarrow$\ locked QC targets\\[5pt]
\textbf{Archived CAMS runs}\ $\longrightarrow$\ unit/time audit\\[5pt]
$\Downarrow$\\[3pt]
\textbf{20:00 IST feature snapshot and availability guard}\\[5pt]
$\Downarrow$\\[3pt]
\textbf{Direct/residual models and matched baselines}*\\[5pt]
$\Downarrow$\\[3pt]
\textbf{Rolling/spatial validation; calibration; advisories}*\\[5pt]
$\Downarrow$\\[3pt]
\textbf{Committed selection lock; fresh final-test approval}*
\end{minipage}}
\caption{Study architecture. Asterisks mark unfinished primary experiment stages. Measured outcomes and forecast fields have separate provenance; final evaluation remains sealed.}
\end{figure}
\subsection{Scope of the contribution}
For station $s$ and target date $d$, the predictor $X_{s,d}$ is formed
before the declared issue time, while $y_{s,d}$ is constructed later from
quality-controlled measurements. The forecast task is next-day prediction,
not concurrent regression on target-day pollutant observations. A high-quality
interface alone cannot establish correction skill. This paper therefore
binds each proposed claim to an immutable result or an explicit pending gate.
The released operational product remains a raw CAMS estimate; offline
station-correction development is a separate research stream.
% END ARCHITECTURE


\clearpage
\section{Data, Targets and Quality}
OpenAQ preserves upstream CPCB-provider station identity; its measurements are
not an official CPCB bulletin or licensed CCR export \cite{openaq2026}.
The inherited catalogue contains 481 eligible stations across 255 parsed cities;
the metadata-only stratified sample selects 181 stations in 100 cities. Current
catalogue selection can introduce survivorship bias and is disclosed.
The locked v3.1 window is 2025-02-19 through 2026-10-07. The protocol was
locked at 2026-10-11T18:32:05.939254+00:00; its prior version and hash remain in the lock history.
Its immutable hash is recorded in the result artifact. This acquisition checkpoint has
179 usable stations, 100 cities, 64799 valid targets and 42481
retained rejected station-days. These are partial-collection counts, not
national study acceptance.

The measured semantic dataset and preliminary input share SHA-256 prefix \texttt{4761a5da60dde059}; the immutable preliminary run has prefix \texttt{287c87db04a6b294}. Full hashes, denominators and artifact locations are recorded in the claim-evidence manifest.

The single CPCB-style engine applies pollutant breakpoints, at least sixteen
distinct hours per pollutant, at least three supported pollutants and a required
particulate species \cite{cpcb2014aqi}. Particulates and the designated gases use
24-hour means; CO and ozone use the maximum complete trailing eight-hour mean.
Targets are labelled AirGuard CPCB-style AQI. Exact zeros, provider flags, poor
coverage and six-hour identical runs are rejected. PM2.5 above PM10 and potential
spikes are flagged and retained. Raw caches and row-addressable rejection ledgers
preserve exclusions. Unsupported ppb conversion is refused.

Severe AQI is an open lower bound of at least 401. The checkpoint contains
383 censored Severe days. They cannot supply exact residual labels or
identified MAE, RMSE, bias or coefficient-of-determination values. Category,
exceedance, forecasts below 401 and lower-bound shortfall remain estimable.
% BEGIN TARGET_METHOD
\subsection{Concentration-to-index transformation}
Within a bounded CPCB interval $[C_L,C_H]$ with index endpoints $[I_L,I_H]$,
the pollutant sub-index is
\begin{equation}
I_p(C)=I_L+\frac{I_H-I_L}{C_H-C_L}(C-C_L),\quad I=\max_{p\in S}I_p.
\end{equation}
The daily support $S$ requires at least three sufficiently covered species,
including particulate matter. It is recorded as an outcome annotation, never
used to select predictors retrospectively. In the current valid development
join, $S$ equals the prerequest frozen PM2.5/PM10/O3 support. Rejected rows
retain reasons and provenance rather than disappearing from denominator
counts. An open concentration interval does not justify linear extrapolation
to a fictitious exact Severe target.
\subsection{Quality-control interpretation}
The six-hour stuck-run rule rejects the complete identical sequence, not only
its final bin. Exact-zero rejection cannot distinguish a genuinely tiny
ozone concentration from a telemetry zero and is therefore a disclosed
preregistered trade-off. Conversely, automatic spike deletion could erase
real pollution episodes. PM2.5/PM10 inconsistencies and potential spikes
remain flagged observations. Sensor rejection, unavailable target and
model error are distinct outcomes; their counts are reported separately.
% END TARGET_METHOD
\section{Issue-Time Availability}
Issue time is 20:00 IST on the preceding day. Observation bins must end by
17:00 IST; daily history is D-2 or earlier. Historical OpenAQ publication time
is unknown, so measured-history analyses explicitly assume a healthy feed.
Archived weather uses the documented 48-hour Previous Runs offset
at \url{https://open-meteo.com/en/docs/previous-runs-api}. The Air Quality API documentation is a separate
source \cite{openmeteo2026}. FIRMS acquisition time alone does not establish release
time. CAMS past fields without their producing run are excluded from the primary
experiment. Every admitted snapshot passes its guard.
\clearpage
\section{Primary Forecast and Correction}
The genuine CAMS global forecast archive \cite{cams2024forecasts} must retain
initialisation, lead, valid time, units and evidence of pre-issue availability.
Reanalysis is prohibited as a replacement. The locked v3.1 protocol specifies
target-support-matched CAMS scoring and a six-pollutant secondary comparator.
Later target support is an evaluation annotation and cannot become a predictor.
The partial join contains 7,560 station-days, 5,506 valid targets and 2,054 retained rejected targets across 42 dates. All 7,560 rows pass the primary availability guard. Every valid target in this subset uses the frozen PM2.5/PM10/O3 support; outcome support remains an evaluation annotation and is excluded from predictors. Full feature integration and the primary runner remain incomplete. The locked comparison is preserved. No primary model has been fitted and no final-test outcomes have been accessed.
All methods share observations and evaluation windows. With frozen CAMS index
$b_{s,d}$ and exact measured index $y_{s,d}$, residual learning uses
\begin{equation}r_{s,d}=y_{s,d}-b_{s,d},\qquad
\widehat y_{s,d}=b_{s,d}+\widehat r_{s,d}.\end{equation}
The ladder comprises raw CAMS, direct ML, residual Ridge, residual LightGBM and
residual XGBoost, with persistence and prior-only climatology as operational
floors. There is no assumption that correction wins. The legacy city model has
a different period, target and input representation, and is reported separately.

\begin{table}[h]\centering\caption{Primary comparison contract}
\begin{tabular}{ll}\toprule Method & Result \\\midrule
Raw operational CAMS & Pending archive \\
Direct ML & Pending archive \\
Residual Ridge & Pending archive \\
Residual LightGBM & Pending archive \\
Residual XGBoost & Pending archive \\\bottomrule\end{tabular}\end{table}






% BEGIN CURRENT CAMS EVIDENCE
\subsection{Verified operational archive checkpoint}
The single accepted probe is a genuine forecast (GRIB type fc), initialized
2025-02-18 00:00 UTC with a 24-hour lead and valid at 2025-02-19 00:00 UTC.
Its 134-byte PM2.5 field contains nine finite values on a 0.4-degree grid.
The development manifest currently verifies 6 of 36 chunks.
Requests cover only the frozen station bounding box and development dates.
Initialization plus lead reconstructs valid time in every admitted field.
Surface PM mass concentrations use $10^9$ to convert kg m$^{-3}$ to
$\mu$g m$^{-3}$. Gas mixing ratios use the same-run lowest-level
temperature $T$, specific humidity $q$ and hybrid pressure $p$:
\begin{equation}
\rho=\frac{p}{287.05T(1+0.6078q)},\quad C_{\mu g/m^3}=10^9\rho x.
\end{equation}
CO instead uses $10^6$ for mg m$^{-3}$. Model-level fields are linearly
interpolated within the same producing run before density conversion;
no extrapolation, reanalysis or post-issue run is admitted. Seven preceding
hourly bins support trailing eight-hour CO/O3 windows.
Historical per-run delivery logs are unavailable: pre-issue availability
uses the documented nominal dissemination bound, explicitly an assumption.
The later ADS retrieval timestamp is provenance, not historical publication.
% END CURRENT CAMS EVIDENCE

\section{Selection and Ablation}
Model settings use training-only inner validation. The locked protocol fixes
90 training days initially, 30 calibration days, 30 test days, a 30-day step and
one-day embargo gaps, with expanding training windows. Each fold is reported.
The latest 90 dates, 2026-07-10 through 2026-10-07, are sealed.
The referenced artifact declares no final-test scoring. The primary selection rule is lowest mean rolling-fold MAE, breaking ties toward Ridge and residual formulations; raw CAMS is retained if no learned method beats it. A committed selection lock, complete isolation audit and fresh user approval are required before final evaluation. Feature ablation separates CAMS, measured history,
forecast weather, calendar, spatial context and optional fires. Fire admission
requires a date-block interval supporting its benefit. An additional sliding
90-day non-CAMS sensitivity is labelled separately.
[RESULT PENDING CAMS ARCHIVE EXPERIMENT]
\clearpage
\section{Temporal and Spatial Verification}
Primary evaluation will report exact-target MAE, RMSE, signed forecast bias,
coefficient of determination, six-category diagnostics, threshold detection,
false alarms and censored-Severe shortfall. Poor, Very Poor and Severe events
receive explicit diagnostics. Seasonal and regional groups below thirty pairs
are listed with counts and withheld metrics.
Spatial tests isolate cities, 25-km connected urban clusters, stations and
regions. Local-history and monitor-free inference are distinct conditions.
Held-out station targets cannot determine training eligibility. The all-station
reference is scored on the identical held-out test rows.

Paired comparison uses seven-day moving date blocks with 5,000 resamples and
a two-sided Diebold--Mariano daily loss comparison with Newey--West variance
and Holm multiplicity adjustment \cite{diebold1995comparing}. A confidence
interval containing zero does not establish improvement.
\section{Preliminary Results Contract}
The separate non-CAMS run completed 24 temporal evaluations,
36 spatial holdouts and 12 calibration-window evaluations;
0 temporal and 0 spatial entries lacked sufficient data.
Its predictions and per-fold metrics come from the exact immutable artifact
recorded in the claim-evidence package. Strict Ridge had lower paired MAE than climatology (signed difference -8.96 AQI; 95\% date-block interval [-12.49, -5.61]; n=38382 exact pairs); lower MAE supported by this interval. History Ridge had higher paired MAE than persistence (signed difference +0.61 AQI; 95\% date-block interval [-0.41, +1.70]; n=31622 exact pairs); a nonzero MAE difference is not established by this interval.
These preliminary comparisons do not establish production approval. They do not answer
whether CAMS station correction improves a genuine operational baseline.
The definitive results table and skill figure remain:
[RESULT PENDING CAMS ARCHIVE EXPERIMENT]

% BEGIN PRELIMINARY FIGURE
\begin{figure}[h]\centering
\setlength{\unitlength}{1mm}
\begin{picture}(78,38)
\put(8,5){\line(1,0){64}}
\put(56,5){\line(0,1){26}}
\put(8.00,4){\line(0,1){2}}
\put(8.00,1){\makebox(0,0){\scriptsize -15}}
\put(24.00,4){\line(0,1){2}}
\put(24.00,1){\makebox(0,0){\scriptsize -10}}
\put(40.00,4){\line(0,1){2}}
\put(40.00,1){\makebox(0,0){\scriptsize -5}}
\put(56.00,4){\line(0,1){2}}
\put(56.00,1){\makebox(0,0){\scriptsize 0}}
\put(72.00,4){\line(0,1){2}}
\put(72.00,1){\makebox(0,0){\scriptsize 5}}
\put(16.04,26){\line(1,0){22.01}}
\put(27.33,26){\circle*{1.5}}
\put(40,31){\makebox(0,0){\scriptsize Strict Ridge minus climatology}}
\put(54.69,13){\line(1,0){6.76}}
\put(57.97,13){\circle*{1.5}}
\put(40,18){\makebox(0,0){\scriptsize History Ridge minus persistence}}
\end{picture}
\caption{NON-CAMS preliminary paired MAE differences and 95\% seven-day date-block intervals. Zero means equal MAE; lower is better. Strict $n=38,382$; history $n=31,622$.}
\end{figure}
% END PRELIMINARY FIGURE
% BEGIN PRELIM_SCOPE
\subsection{Limits of the preliminary experiment}
The non-CAMS plan uses 30 valid training days per station and a fixed Ridge
penalty of 10; these choices must not be mistaken for primary CAMS tuning.
The locked primary station rule remains 120 valid training days. Thus the
first 90-day primary training window cannot meet station eligibility; it
must be reported as insufficient, not rescued by future coverage or a
silent threshold reduction. Later windows are evaluated when sufficient.
Primary model selection has not begun. Absence of a primary score is not
evidence of equivalence, and preliminary skill against climatology cannot
establish benefit over archived CAMS or persistence.
% END PRELIM_SCOPE
\section{Reproducibility}
Files and decompressed cache content have SHA-256 manifests. Dataset identity,
feature version, protocol hash, source hash and run configuration accompany each
result. Local station reconstruction is audited independently of derived caches.
A clean-clone measured-dataset reconstruction passed its semantic hash audit under the recorded runtime. A complete primary reproduction remains pending. Numerical values
in this skeleton are generated from evidence artifacts, never filled manually.
\clearpage
\section{Uncertainty Calibration}
Split conformal intervals at 80\% and 90\%, quantile calibration and
forecast-category Mondrian calibration use only the preceding calibration
window. Coverage and widths are reported by category, season and region.
For a censored Severe outcome, an upper bound reaching 401 is merely compatible
with the known lower bound; it does not identify true-value coverage.
Calibration under temporal and spatial shift is evaluated empirically;
exchangeability guarantees are not asserted for dependent station-days
\cite{gibbs2021aci,xu2021enbpi}. Proper uncertainty scoring requires identified
outcomes \cite{gneiting2007scoring}.





% BEGIN PRELIMINARY COVERAGE
\begin{table}[h]\centering\caption{NON-CAMS preliminary first-fold calibration failures}
\begin{tabular}{lrr}\toprule Method & Exact coverage & Mean width\\\midrule
80\% split & 0.855 & 95.5\\
90\% split & 0.944 & 120.6\\
\bottomrule\end{tabular}
\end{table}
Marginal coverage conceals important failures: neither interval covered any
of the 35 exact Very Poor outcomes in this first fold. Poor-or-worse coverage
was also zero on 48 exact pairs. These are preliminary validation failures,
not primary CAMS or final-test findings; they preclude claiming reliable
high-pollution uncertainty from average coverage alone.
% END PRELIMINARY COVERAGE

\section{Vulnerability-Aware Advisory Evaluation}
Declared exposure strata include general adults, children, older adults,
asthma/COPD, heart conditions, pregnancy and outdoor workers. They are not
personal medical-risk predictions. P0 uses the point forecast; P1 the calibrated
median; P2 the calibrated 0.75 quantile; P3 the smallest calibration-window
candidate achieving the corrected empirical under-warning aim $(k+1)/(n+1)$.
The aims are 10\% for general adults and 5\% for the other declared strata.
If no candidate qualifies, the policy reports that fact.
Under-warning, over-warning, serious under-warning and high-pollution misses
are reported alongside date-block confidence intervals. A lower warning rate
does not demonstrate a health benefit. Unsupported cohort coefficients remain
unadopted. [RESULT PENDING CAMS ARCHIVE EXPERIMENT]
% BEGIN ADVISORY_METHOD
\subsection{Policy losses and calibration boundary}
For ordered observed category $c$ and warning category $w$, the basic losses
are $L_-=\mathbf{1}\{w<c\}$ and $L_+=\mathbf{1}\{w>c\}$.
Serious under-warning, missed Poor-or-worse and missed Very-Poor-or-worse
events remain separate diagnostics. Lower under-warning can require more
false alarms; neither objective is optimized on a test window.
The candidate quantiles are fixed before evaluation. P3 chooses on the
calibration window only and then reports its frozen test-window behavior,
including failure to find any feasible candidate. Its corrected empirical
rate $(k+1)/(n+1)$ is a policy selection statistic, not an individual disease
probability or a proven dependent-data risk bound.
\subsection{Finite-sample intervals and failure reporting}
For exact calibration residuals $e_i=|y_i-\widehat y_i|$, split conformal
uses the order statistic with rank $\lceil(n+1)(1-\alpha)\rceil$.
The interval is $[\max(0,\widehat y-q),\widehat y+q]$; missing or inadequate
calibration support is reported rather than assigned a radius of zero.
Mondrian groups use the forecast category, so observed test categories cannot
choose an interval. Quantile crossing is repaired using predictions only.
Calibration labels exclude censored Severe values from exact residuals;
Severe category/threshold behavior is still reported. Under temporal/spatial
dependence, coverage is an empirical finding, not automatically a theorem.
\subsection{Exposure guidance and clinical scope}
Declared vulnerability strata set environmental-warning tolerances. Shared
forecast outcomes do not constitute separate clinical cohort studies. The
system does not predict diagnoses, prescribe medication, estimate personal
health probabilities or demonstrate avoided illness. Its interface labels
measurement versus modelling, observation age, unsupported intervals and
raw-model uncertainty. Clinical coefficients remain unvalidated.
% END ADVISORY_METHOD
\section{Prospective Verification}
The current audited ledger contains 24 real forecasts across 2 target days, and 0 verified days. 0 forecasts have a bound station. 12 historical forecasts lack an issuance Git SHA; missing historical identities will not be invented. New issuance now records its deployment identity. Unbound records cannot be made station forecasts by retrospectively selecting an outcome monitor. No experimental correction is deployed. Issued forecasts remain immutable; a thirty-day track record requires real elapsed days and matched measured outcomes. Station-level prospective evaluation remains blocked until stations are bound before issuance.
\clearpage
\section{Limitations and Planned Completion}
ADS setup is complete. Incomplete archive coverage and primary experiments block the genuine CAMS comparison. Partial station acquisition, current-catalogue survivorship,
unknown historical observation latency and limited extreme-event support bound
the preliminary analysis. Optional fire release time remains unproven. The
initial station-free prospective forecasts do not establish station forecast
skill. Final-test isolation, transparent artifacts and explicit missing results
are necessary conditions for completion, not evidence of improved accuracy.
The intended completed narrative is CAMS baseline, India-specific correction,
rolling/spatial testing, calibrated uncertainty, cohort advisory trade-offs and
prospective station verification.
\section{Conclusion}
The study preserves its locked protocol, measured targets and immutable preliminary experiment; genuine CAMS alignment is now audited on a partial development subset. Correction performance,
operational intervals, advisory risk attainment and thirty-day verification
remain hypotheses. [RESULT PENDING CAMS ARCHIVE EXPERIMENT]
\section*{Ethics and AI Disclosure}
No human-subject data was collected. Cohort guidance is environmental exposure
advice, not validated medical advice. Any participant study requires prior
ethics approval and consent. Claude Code and OpenAI Codex assisted software,
tests and drafting. Authors must review the methods, bibliography, disclosures
and generated claims before submission. The named author must approve the final manuscript before submission.
\begin{thebibliography}{99}
\bibitem{jena2021delhi} Jena, Chinmay and Ghude, Sachin D. and Kumar, Rajesh and Debnath, Sreyashi and Govardhan, Gaurav and Soni, Vijay K. and Kulkarni, Santosh H. and Beig, G. and Nanjundiah, Ravi S. and Rajeevan, M., ``Performance of high resolution (400 m) PM2.5 forecast over Delhi'', Scientific Reports, Springer Science and Business Media LLC, vol. 11, no. 1, 2021, DOI: \url{https://doi.org/10.1038/s41598-021-83467-8}, \url{http://dx.doi.org/10.1038/s41598-021-83467-8}.
\bibitem{diebold1995comparing} Diebold, Francis X. and Mariano, Roberto S., ``Comparing Predictive Accuracy'', Journal of Business \& Economic Statistics, Informa UK Limited, vol. 13, no. 3, pp. 253--263, 1995, DOI: \url{https://doi.org/10.1080/07350015.1995.10524599}, \url{http://dx.doi.org/10.1080/07350015.1995.10524599}.
\bibitem{gneiting2007scoring} Gneiting, Tilmann and Raftery, Adrian E, ``Strictly Proper Scoring Rules, Prediction, and Estimation'', Journal of the American Statistical Association, Informa UK Limited, vol. 102, no. 477, pp. 359--378, 2007, DOI: \url{https://doi.org/10.1198/016214506000001437}, \url{http://dx.doi.org/10.1198/016214506000001437}.
\bibitem{xu2021enbpi} Chen Xu and Yao Xie, ``Conformal prediction for time series'', 2020, arXiv: \url{https://arxiv.org/abs/2010.09107}.
\bibitem{gibbs2021aci} Isaac Gibbs and Emmanuel Candès, ``Adaptive Conformal Inference Under Distribution Shift'', 2021, arXiv: \url{https://arxiv.org/abs/2106.00170}.
\bibitem{cpcb2014aqi} {Central Pollution Control Board}, ``National Air Quality Index'', Central Pollution Control Board, Ministry of Environment, Forest and Climate Change, India, Control of Urban Pollution Series CUPS/82/2014-15, 2014, Online: \url{https://cpcb.nic.in/National-Air-Quality-Index/}.
\bibitem{cams2024forecasts} {Copernicus Atmosphere Monitoring Service}, ``CAMS global atmospheric composition forecasts'', Copernicus Atmosphere Data Store, 2024, Online: \url{https://ads.atmosphere.copernicus.eu/datasets/cams-global-atmospheric-composition-forecasts}.
\bibitem{openaq2026} {OpenAQ}, ``OpenAQ API version 3 documentation'', OpenAQ, 2026, Online: \url{https://docs.openaq.org/}.
\bibitem{openmeteo2026} {Open-Meteo}, ``Air Quality API documentation'', Open-Meteo, 2026, Online: \url{https://open-meteo.com/en/docs/air-quality-api}.
\bibitem{riccio2024mos} A. Riccio and E. Chianese, ``Technical note: Accurate, reliable, and high-resolution air quality predictions by improving the Copernicus Atmosphere Monitoring Service using a novel statistical post-processing method,'' Atmos. Chem. Phys., vol. 24, pp. 1673--1689, 2024, doi: 10.5194/acp-24-1673-2024.
\bibitem{petetin2022mos} H. Petetin et al., ``Model output statistics (MOS) applied to Copernicus Atmospheric Monitoring Service (CAMS) O$_3$ forecasts: trade-offs between continuous and categorical skill scores,'' Atmos. Chem. Phys., vol. 22, pp. 11603--11630, 2022, doi: 10.5194/acp-22-11603-2022.
\bibitem{colette2025cams} A. Colette et al., ``Copernicus Atmosphere Monitoring Service--Regional Air Quality Production System v1.0,'' Geosci. Model Dev., vol. 18, pp. 6835--6883, 2025, doi: 10.5194/gmd-18-6835-2025.
\end{thebibliography}
\end{document}
